% LaTeX source for an arXiv-style conceptual paper on "Empty Mathematics (Ku-Mathematics)"
\documentclass[11pt]{article}
\usepackage{amsmath, amssymb, amsthm, amsfonts}
\usepackage{hyperref}
\usepackage{fullpage}
\usepackage{times}

\title{Toward a Theory of ``Empty Mathematics'':\\ A Foundational Framework for Observer-Dependent Abstraction}
\author{Masaomi Chiba\\ independent researcher}
\date{\today}

\begin{document}
\maketitle

\begin{abstract}
This paper introduces a conceptual framework called \emph{Empty Mathematics} (``Ku-Mathematics''),
intended as a pre-foundational layer underlying category theory, model theory, and other modern abstractions.
Empty Mathematics formalizes two key cognitive operations:
(i) the \emph{emptying} (``ku-ization'') of structures into observer-dependent coherent states, and
(ii) the \emph{coloring} (``shiki-ization'') of these states into determinate mathematical objects.
We argue that this bidirectional process provides a more primitive and flexible substrate than traditional foundations,
enabling a principled account of multi-universe reasoning, observer variability, and non-isomorphic translations.
We further propose that this framework is required to model certain extreme mathematical theories,
such as Inter-Universal Teichmüller Theory (IUT), which inherently rely on incompatible or non-translatable structures.
\end{abstract}

\section{Introduction}
Classical foundations of mathematics---set theory, type theory, and category theory---presuppose that
mathematical structures exist in a single determinate universe equipped with stable logical rules.
However, several advanced theories (e.g., anabelian geometry, relative motives, and IUT) expose phenomena
that cannot be internalized in a single universe without destructive loss of structure.

We introduce \emph{Empty Mathematics} as a meta-level framework that explicitly models
(1) observer dependence, (2) multi-universe inconsistency, and
(3) non-resolution of certain structural relations.
This framework extends beyond category theory by incorporating non-translatable states and
systematically handling the alternation between undetermined and determined mathematical states.

\section{Coherent and Decoherent Mathematical States}
We define two complementary modes:

\begin{itemize}
    \item \textbf{Coherent (Empty) States}: mathematical entities remain in an indeterminate state;
    morphisms, objects, and even universes are not yet fixed. These states represent ``pre-mathematical''
    abstraction.
    \item \textbf{Decoherent (Colored) States}: specific mathematical structures become fixed;
    categories, sets, elements, and morphisms acquire determinate identity.
\end{itemize}

Empty Mathematics studies the controlled alternation between these states. Category theory corresponds to a fully decoherent phase, where composition, identity, and naturality are strictly enforced. In contrast, coherent states allow incompatible identifications and alternative observer positions.

\section{Observer-Dependent Universes and F-Fix Points}
Let $F$ denote an observer configuration (``F-fix point''). A coherent state is defined relative to $F$.
Different observers may induce incompatible decoherent states even when starting from identical coherent seeds.

This formalizes the notion that mathematical truth, when abstract enough, becomes partially observer-relative.

\section{Non-Isomorphic Translations}
One of the motivations for Empty Mathematics is to give a principled account of ``non-isomorphic translations''---transitions between mathematical universes where no natural functor or equivalence exists.
Empty Mathematics explicitly allows ``colorings'' that cannot be compared within any category.

Structures such as those encountered in IUT---where certain identifications are intentionally destroyed---fit naturally into this framework.

\section{Toward Applications}
Potential applications include:

\begin{itemize}
    \item a meta-logic for multi-universe mathematics,
    \item a foundation for future AI systems capable of observer-variable reasoning,
    \item a conceptual simplification of extreme theories such as IUT,
    \item a cognitive model explaining creativity in high-level mathematics.
\end{itemize}

\section{Conclusion}
Empty Mathematics proposes a pre-foundational abstraction layer that models both coherent and decoherent phases of mathematical reasoning. Future work includes formalizing observer shifts, defining coherent-to-decoherent operators, and exploring connections with homotopy type theory and higher category theory.

\end{document}
